\begin{helper}[Step 1 zero-divider non-PE branch]
Let $T$ be a positive-rung ETT of rung $n$ on a finite loopless
multigraph with individual edge identities, generating coloring $\psi$,
maximum property, properness, and exactly the distinguished edge $e_0$
uncolored.  Assume uniformly for every positive-rung maximum ETT of rung
strictly below $n$ the complete six assertions used in Step 1:
terminal elementarity; stable interchangeability of the final ladder
closure, including isolated path components; PE stability descent unconditionally before each support and, when the
first-appearance index is greater than one, to the preceding ladder stage;
stored pairs at distinct supports with the same first appearance are disjoint; the terminal PE component intersection assertion; the
defective-endpoint precedence assertion; and the terminal mod-coloring
assertion.  Assume that the terminal extension kind is not PE.

Let $B$ be a hierarchy base and $H$ a hierarchy over $B$ with zero
dividers.  Put $A=H.\mathit{starred}(0)$, the exact hierarchy core.  Let
$a\ne b$ be colors such that $a$ is missing and singleton-closed on $A$.
Suppose there are exact $(A,\psi,\{a,b\})$-exit paths $P_1,P_2$ with
distinct outside endpoints and disjoint incidence-walk supports.  Then
these data are contradictory.
\end{helper}
\begin{proof}
Write $n$ for the rung number, $T_n$ for the tree in series slot $n$,
$\psi_{n-1}$ for the coloring in slot $n-1$, and $D_i$ for the exact
computed stored-color set in slot $i$.  Thus the predecessor uses
$D_{n-1}$, while the terminal PE assertions use $D_n$.  The
generating-coloring field of \noderef{ExtendedTashkinovTree} identifies
the displayed coloring $\psi$ with series color slot $n$.  Since
the divider count is zero, the value of
\noderef{HierarchyStarredStage} at its only index is the core stored in
the given \noderef{HierarchyBase}.  Thus the tree called $A$ in the
statement is exactly that core.

We first record two consequences of each alleged exit path.  Unfold
\noderef{TreeExitPath}.  Its node list starts at the outside vertex-node
$a_i$, ends at the tree vertex-node $b_i$, has length at least two, is
simple, and contains no other vertex-node of $A$.  By the bipartite
adjacency definition in \noderef{TwoColorIncidenceGraph}, the node just
before $b_i$ is an exact edge identity $f_i$ colored $a$ or $b$ and
incident with $b_i$.  The vertex-node before $f_i$ is outside $A$,
whereas $b_i$ is inside.  Hence the endpoint description in
\noderef{BoundaryEdges} gives $f_i\in\partial_G(A)$, and
\noderef{BoundaryColorEdges} gives
\[
 f_i\in\partial_{\psi,a}(A)\quad\hbox{or}\quad
 f_i\in\partial_{\psi,b}(A).                         \tag{1}
\]
Similarly, the node after $a_i$ is an exact incident edge $g_i$ colored
$a$ or $b$.  This first edge will certify below that the outside endpoint
is not isolated after the recoloring.  These conclusions use actual edge
nodes, so parallel edge identities are not identified.

By the terminal-kind hypothesis we are in the non-PE case.
The nonparallel field of \noderef{HierarchyBase} identifies
$A=T_n$.  We claim that
\[
 d\in\overline\psi(T_n)\quad\Longrightarrow\quad
 \partial_{\psi,d}(T_n)=\varnothing                 \tag{2}
\]
for every color $d$.  To prove it, choose, using
\noderef{TreeMissingColors}, a vertex $y\in V(T_n)$ with
$d\in\overline\psi(y)$.  Positivity of $n$ lets us inspect the final
certificate supplied by \noderef{TashkinovSeries}.  It cannot be the
initial constructor of \noderef{TashkinovSeriesStep}, and the present
case excludes its PE constructor, so it is RE or SE.

In the RE case, the certificate says both that $V(T_n)$ is closed under
$\psi_{n-1}$ and that $\psi=\psi_{n-1}$.  If an identity $f$ belonged to
$\partial_{\psi,d}(T_n)$, then
\noderef{BoundaryColorEdges} would give
$f\in\partial_G(T_n)$ and $\psi(f)=d$.  Together with the missing witness
at $y$, this contradicts the quantified condition in
\noderef{ColoredVertexSetClosed}.  Hence no such $f$ exists.

In the SE case, its \noderef{MaximumDefectiveChoice} has selected
coloring $\psi$, and its selected-stability field says that $\psi$ is
stable on $T_n$ relative to $\psi_{n-1}$.  Both colorings are proper and
leave exactly $e_0$ uncolored by the same two certificates.  Symmetry in
\noderef{sc2} therefore says that $\psi_{n-1}$ is stable relative to
$\psi$.  The vertex clause of \noderef{StableColoring} transports
$d\in\overline\psi(y)$ to
$d\in\overline{\psi_{n-1}}(y)$.  If
$f\in\partial_{\psi,d}(T_n)$, then $f$ is incident with $T_n$, and the
incident-edge clause of that reverse stability applies because $d$ is
tree-missing under $\psi$; it yields
$\psi_{n-1}(f)=\psi(f)=d$.  The preceding-coloring closedness field of
the SE certificate now contradicts \noderef{ColoredVertexSetClosed} at
$f,d,y$.  This proves (2) in both possible constructors.  Applying the
definition \noderef{TreeColorSetClosed} to (2) says, in particular, that
every singleton $\{d\}$ with
$d\in\overline\psi(T_n)$ is closed on $T_n$.

The color $b$ is not tree-missing on $T_n$.  Otherwise (2) makes
$\partial_{\psi,b}(T_n)$ empty, while the assumed singleton closedness
of $a$ makes $\partial_{\psi,a}(T_n)$ empty.  This contradicts the
exact boundary identity $f_1$ in (1).  By
\noderef{TreeMissingColors}, $b$ is therefore missing at no vertex of
$T_n$; by the complement definition \noderef{MissingColors}, it belongs
to \noderef{PresentColors} at every such vertex.

Let $x$ be the root vertex of the nonempty tree-sequence $T_n$, as
provided by \noderef{TreeSequence} and
\noderef{TreeSequenceVertexSet}.  We now construct a color missing at
$x$, rather than merely asserting one exists.  At each
$d\in\operatorname{PresentColors}(\psi,x)$, choose an exact edge identity
$h_d$ incident with $x$ in the sense of \noderef{Incident} and colored
$d$; this is the witness in \noderef{PresentColors}.  Properness makes
the incident edge of a fixed present color unique, and distinct colors
have distinct chosen edge identities, since the value $\psi(h_d)$ cannot
be two different colors.
Thus $d\mapsto h_d$ injects the finite set of present colors into the
finite set of edge identities incident with $x$.  The latter set has
cardinality $\deg_G(x)$ by \noderef{Degree}, and
$\deg_G(x)\leq\Delta(G)$ by \noderef{MaximumDegree}.  The certified
initial constructor in \noderef{TashkinovSeries} supplies the initial
\noderef{KTriple}, whose degree clause is
$\Delta(G)+1\leq k$.  Consequently
\[
 |\operatorname{PresentColors}(\psi,x)|
 \leq\deg_G(x)\leq\Delta(G)<k=|[k]|.
\]
Choose $c\in[k]$ outside the present-color set.  By
\noderef{MissingColors}, $c\in\overline\psi(x)$, so
$c\in\overline\psi(T_n)$ by \noderef{TreeMissingColors}; equation (2)
makes $\{c\}$ closed on $T_n$.  On the other hand, the preceding
paragraph makes $b$ present at $x$.  A color cannot be both present and
missing under the complementary definitions of \noderef{PresentColors}
and \noderef{MissingColors}, so
\[
                         c\ne b.                  \tag{3}
\]

We next build the changed coloring and check both exit paths.  This
requires a separate equality case.

If $a=c$, set $\rho=\psi$.  It is proper and leaves exactly $e_0$
uncolored by the hypotheses, with the latter predicate understood as in
\noderef{IsOneEdgeUncolored}.  Reflexivity in \noderef{sc2} makes it
$(T_n,C,\psi)$-stable for every control set $C$.  Replacing the displayed
name $a$ by its equal name $c$ in the two given walks and their
\noderef{TreeExitPath} certificates gives two $(c,b)$ exit paths with
the same supports.

Suppose now that $a\ne c$.  Define the following finite root set in the
$(a,c)$ incidence graph:
\[
 \mathcal R=\{\operatorname{inl}a_1,\operatorname{inl}a_2\}
 \mathbin\cup
 \{\operatorname{inr}f: \operatorname{inr}f\in
       \operatorname{supp}P_1\cup\operatorname{supp}P_2,
       \ \psi(f)=a\}.                              \tag{4}
\]
It is finite because each walk support is a finite list.  Every full
$(a,c)$ component rooted in $\mathcal R$ avoids $V(T_n)$.  Here are the
details.  A root $\operatorname{inl}a_i$ is already outside $T_n$ by
\noderef{TreeExitPath}.  For a root $\operatorname{inr}f$ in the second
set, $f$ lies on one of the exit walks and has color $a$.  It cannot be
the final edge $f_i$, since that edge is a boundary edge and
$\{a\}$ is closed.  Hence one of its adjacent vertex-nodes on that walk
occurs before the unique tree vertex $b_i$ and is outside $T_n$.  In
either case the rooted component contains an outside vertex-node.

If such a component also contained a vertex-node of $T_n$, take a
shortest incidence path inside the component from an outside vertex-node
to a vertex-node of $T_n$, and stop at its first vertex-node in $T_n$.
The preceding vertex-node is outside $T_n$, and between them is one exact
edge-node $h$.  By \noderef{TwoColorIncidenceGraph}, $h$ is incident
with both vertices and $\psi(h)=a$ or $\psi(h)=c$.
Thus \noderef{BoundaryEdges} puts $h$ in $\partial_G(T_n)$ and
\noderef{BoundaryColorEdges} puts it in
$\partial_{\psi,a}(T_n)$ or $\partial_{\psi,c}(T_n)$.  Both alternatives
are impossible by the singleton closedness of $a$ and $c$.  This proves
the required outside-component condition for every root in (4).

Let $\rho$ be the coloring which swaps $a,c$ on exactly the union of
these rooted full components and fixes every other edge identity.  The
distinctness $a\ne c$, the root set (4), the preceding outside proof, and
the edge-by-edge displayed rule are exactly
\noderef{OutsideKempeRecoloring}.  Therefore
\noderef{OutsideKempeRecoloringStable} says that $\rho$ is proper, leaves
exactly $e_0$ uncolored, is $(T_n,C,\psi)$-stable for every $C$, and
fixes every exact edge identity incident with $T_n$.

For $i=1,2$, reinterpret the same finite node list as a walk $P_i'$ in
$I_G(\rho;c,b)$.  To check adjacency, consider an edge-node $f$ of the
old walk.  By \noderef{TwoColorIncidenceGraph}, its old color is $a$ or
$b$.  If it is $a$, then $\operatorname{inr}f$ was put in (4), so the
defining formula of \noderef{OutsideKempeRecoloring} gives
$\rho(f)=c$.  If it is $b$, then (3) and the original $a\ne b$ show that
$b$ is neither exchanged color, so $\rho(f)=b$.  Incidence with the
neighboring vertex-nodes is unchanged.  Hence every consecutive pair is
adjacent in the new incidence graph.  The list is unchanged, so it is
still simple, has the same positive length, has endpoints $a_i,b_i$, and
has the same support; in particular its only tree vertex is still $b_i$.

It remains to check the missing-color alternative at $a_i$.  If $b$ was
missing there, every $b$-edge is fixed, so it remains missing.  If $a$
was missing there, the root $\operatorname{inl}a_i$ belongs to (4).
Every $a$- or $c$-edge incident with $a_i$ lies in that full rooted
component by \noderef{TwoColorComponent}, and all of them are swapped.
Thus a $c$-edge is present at $a_i$ under $\rho$ exactly when an $a$-edge
was present there under $\psi$.  Since the latter does not exist, $c$ is
missing under $\rho$.  Using \noderef{PresentColors} and then its
complement \noderef{MissingColors} makes this implication exact.  We
have now checked every field of \noderef{TreeExitPath} for $P_i'$.
Because their node lists did not change, the assumed support disjointness
also persists.  The equality and nonequality constructions therefore
give the same package: a proper, one-uncolored coloring $\rho$ stable on
$T_n$, and two support-disjoint $(T_n,\rho,\{c,b\})$ exit paths.

We now construct the full new path components rather than transporting
old component identities.  Stability at $x$ preserves its missing and
present sets.  Hence $c$ is missing and $b$ is present at $x$ under
$\rho$.  The present-color witness gives one incident $b$-edge.  There is
no incident $c$-edge, and properness as defined in
\noderef{IsProperEdgeColoring} permits at most one incident $b$-edge.
Thus the vertex-node $\operatorname{inl}x$ has exactly one neighbor in
$I_G(\rho;c,b)$.

For each outside endpoint $a_i$, the first edge of $P_i'$ is incident
with $a_i$ and has color $c$ or $b$.  Its exit-path missing alternative
says that one of these colors is missing.  If $c$ is missing, the first
edge has color $b$; if $b$ is missing, it has color $c$.  Properness
excludes a second incident edge of that same color, while missingness
excludes every incident edge of the other color.  Thus each
$\operatorname{inl}a_i$ also has exactly one neighbor.  The general
degree bound in \noderef{TwoColorIncidenceDegreeAtMostTwo} is consistent
with these exact one-neighbor calculations.

Let $K_x,K_1,K_2$ be the full \noderef{TwoColorComponent}s rooted at
$\operatorname{inl}x,\operatorname{inl}a_1,
\operatorname{inl}a_2$, respectively.  Equation (3), membership of each
root in its own component, and the one-neighbor witnesses prove, by the
definition \noderef{TwoColorPathComponent}, that all three are path
components.  The component $K_x$ meets $T_n$ at $x$.  Each $P_i'$ is a
walk from $a_i$ to $b_i\in V(T_n)$, so connectedness in
\noderef{TwoColorComponent} puts $\operatorname{inl}b_i$ in $K_i$;
hence $K_i$ also meets $T_n$.

The three leaf nodes are distinct: $a_1\ne a_2$ is assumed, and $x$ is
in $T_n$ whereas both $a_i$ are outside.  They cannot all lie in one full
component.  Indeed, if $K_x=K_1=K_2$, apply
\noderef{TwoColorComponentSpanningPath} to that component and the proper
coloring $\rho$.  It supplies one simple path whose support is the whole
component.  Any one-neighbor node in the support must be an endpoint:
an internal node of a simple path has two distinct predecessor and
successor neighbors.  The spanning path has only its two endpoints, so
it cannot contain the three distinct one-neighbor nodes $x,a_1,a_2$.
Therefore $K_x\ne K_i$ for at least one $i\in\{1,2\}$.  We have produced
two distinct full $(c,b)$ path components which both meet $T_n$.

It remains to invoke the correct rung of interchangeability.  If $n>1$,
\noderef{CertifiedETTPredecessor} gives a genuine rung-$n-1$ ETT whose
displayed terminal tree is exactly $T_n$, whose generating coloring is
exactly $\psi_{n-1}$, and which has maximum property.  Its pointwise
data equalities identify its final tree slot and coloring slot with the
original ones; crucially, its computed-control equality for
\noderef{TashkinovStoredColors} identifies its final control set with
the original $D_{n-1}$.  In the RE case $\psi=\psi_{n-1}$, so the
stability of $\rho$ already has the required reference coloring.  In the
SE case the selected-stability field makes $\psi$ stable relative to
$\psi_{n-1}$ with control $D_{n-1}$, and transitivity in
\noderef{sc2} composes it with the just-proved stability of $\rho$
relative to $\psi$.  Thus lower assertion (ii) applies because $c$ is
missing at $x\in V(T_n)$.

If $n=1$, the retained tree and coloring are precisely the initial tuple
$T_1,\psi_0$.  The RE equality or the same SE stability composition gives
the exact stable-coloring premise of \noderef{TechnicalTheoremBase}; its
interchangeability conclusion again applies because $c$ is missing on
$T_n$.  In either rung case, unfolding \noderef{TreeInterchangeable}
forces any two $(c,b)$ path components meeting $T_n$ to be the same.
This contradicts the distinct pair constructed above and completes the
non-PE case.

The contradiction above proves the non-PE branch.
\end{proof}
